\chapter{Analyse de l'existant}

\section{Architecture actuelle de la chaîne d'alerting}

L'architecture étudiée repose sur un ensemble de composants AWS permettant de créer, maintenir et diffuser des
alarmes techniques autour des services déployés sur ECS.

Le fonctionnement observé est le suivant :
\begin{itemize}
	\item des changements sur les services (création, mise à jour) sont détectés via EventBridge,
	\item des fonctions Lambda orchestrent la création ou la suppression d'alarmes CloudWatch,
	\item les alarmes déclenchées produisent des notifications (SNS puis canaux de diffusion, par exemple Slack).
\end{itemize}

Cette chaîne automatise la gestion des alarmes mais génère un niveau de complexité élevé, avec de nombreux points
de configuration et de couplage entre services techniques.

\begin{figure}[h]
	\centering
	\fbox{%
		\parbox{0.9\textwidth}{%
			\textbf{Schéma de l'architecture actuelle d'alerting (à insérer).}\\
			Cette figure représentera le flux principal : ECS $\rightarrow$ EventBridge $\rightarrow$ Lambda
			$\rightarrow$ CloudWatch $\rightarrow$ SNS/Slack.
		}%
	}
	\caption{Vue d'ensemble de la chaîne actuelle de création et diffusion d'alarmes}
	\label{fig:archi-alerting-actuelle}
\end{figure}

\section{Constats opérationnels sur la pertinence des alertes}

L'étude de l'existant met en évidence une limite majeure : une partie significative des alertes est construite autour
de métriques de causes (CPU, mémoire, etc.) plutôt que de symptômes directement corrélés à un impact utilisateur.

Cette situation entraîne plusieurs effets :
\begin{itemize}
	\item des réveils d'astreinte pour des situations parfois non critiques,
	\item une fatigue liée au volume de notifications,
	\item une difficulté à identifier rapidement le service réellement en défaut.
\end{itemize}

L'efficacité globale de la réponse à incident est donc dégradée, malgré une instrumentation technique présente.

\section{Problème posé et objectifs de transformation}

Le besoin exprimé est d'introduire des niveaux de criticité explicites dans le système d'alerting, afin de mieux
prioriser les interventions.

Les objectifs fonctionnels retenus sont :
\begin{itemize}
	\item alerter la bonne personne sur le bon service,
	\item limiter les notifications à faible valeur opérationnelle,
	\item améliorer la visibilité sur l'état réel de la plateforme,
	\item réduire les traitements d'urgence subis en favorisant l'anticipation.
\end{itemize}

Cette problématique constitue le fil directeur des chapitres suivants.

\section{Synthèse du diagnostic}

Le diagnostic de départ peut être résumé en trois points :
\begin{itemize}
	\item l'infrastructure d'alerting est riche mais hétérogène,
	\item la pertinence métier des alertes reste insuffisante dans certains cas,
	\item la priorisation des interventions n'est pas encore assez explicite.
\end{itemize}

Ce constat justifie la définition d'un cadre de criticité formalisé, présenté dans le chapitre suivant.